\documentclass[10pt,a4paper]{amsart}
\usepackage[margin=19mm]{geometry}
\usepackage{amsmath,amssymb,mathtools}
\usepackage[scr=rsfs,cal=euler]{mathalfa}
\usepackage{xfrac}
\usepackage[unicode,hidelinks,bookmarks=false]{hyperref}
\newcommand{\ie}{i.e.\ }
\newcommand{\cf}{cf.\ }
\newcommand{\resp}{respectively\ }
\newcommand{\Subsection}{\subsection}
\providecommand{\nobreakdash}{\nobreak\hbox{-}}
\setlength{\emergencystretch}{2em}
\multlinegap0pt
\usepackage{mathtools}
\usepackage{amssymb}
\usepackage{xcolor}
\usepackage{dsfont}
\usepackage[shortlabels]{enumitem}
\newtheorem{definition}{Definition}
\newtheorem{proposition}{Proposition}
\newcommand{\R}{{\mathbb R}}
\title{\textbf{A Dynamical Systems Perspective \\ on the Analysis of Neural Networks}}
\author{Dennis Chemnitz, Maximilian Engel, Christian Kuehn, Sara-Viola Kuntz}
\date{Source: arXiv:2507.05164v2 (CC BY 4.0)}
\begin{document}
\maketitle

\noindent\emph{Source and licence.} \textbf{A Dynamical Systems Perspective \\ on the Analysis of Neural Networks}. Version arXiv:2507.05164v2 (CC BY 4.0), distributed by arXiv under the Creative Commons Attribution 4.0 International licence, http://creativecommons.org/licenses/by/4.0/, as recorded in the arXiv licence metadata for this exact version and shown on the abstract page https://arxiv.org/abs/2507.05164v2. Full licence notice: \texttt{licenses/arxiv-2507.05164-CC-BY-4.0.txt}.\par\smallskip

\noindent\textit{Editorial scope.} Counted unit: the article's proposition cekk\_prop:LocGlobEqui (source0.tex lines 706--708). The article's complete original proof of it is delivered with it (source0.tex lines 709--719, one \texttt{proof} environment of 11 lines). No mathematics is written, repaired or re-derived: the statement and the proof are the article's own lines, in the article's own notation, and no retained line is substituted or re-flowed.  Retained uncounted as the source-local context the proof needs: lines 686--693.  Nothing was omitted from any retained span, so the package carries no omission record.  No retained line contains a citation, so the wrapper carries no bibliography and no bibliography entry.  No retained line contains a figure environment, an image-inclusion command or drawing code, so no image is delivered and the package declares no asset dependency.  Editorial formatting only, in addition: an AMS \texttt{amsart} preamble that restates the article's own declarations and macros used by the retained lines, namely the environments \texttt{definition}, \texttt{proposition} on separate counters, together with the standard packages those lines need.  The retained environments print in this document as Definition 1, Proposition 1 in order of appearance, whereas the article prints the same items with the numbering of its own full document.  The complete native source of the article as distributed in the arXiv e-print for this exact version is bundled unchanged as \texttt{sources/181160/source0.tex}.
\begin{definition}[Milnor Stability of an Isolated Fixed Point]\label{cekk_def:Milnor}~
	\begin{enumerate}
		\item[(i)]\index{Dynamical Stability!for Isolated Minima}An isolated fixed point $\theta^* = \varphi(\theta^*)$ is called \emph{Milnor stable} if the set of possible initializations $\theta_0\in \R^D$ for which the sequence of gradient descent iterates $\theta_n$ converges to $\theta^*$ has a positive Lebesgue measure, i.e., if
		$$\operatorname{Vol}\left\{\theta \in \R^D : \lim_{n \to \infty} \varphi^n(\theta) = \theta^*\right\}>0.$$
		\item[(ii)]An isolated fixed point $\theta^* = \varphi(\theta^*)$ is called \emph{locally Milnor stable} if for every open neighborhood $U\subset \R^D$ of $\theta^*$ the set of possible initializations $\theta_0\in U$, for which the sequence of gradient descent iterates $\theta_n$ stays in $U$ forever and converges to $\theta^*$, has a positive Lebesgue measure, i.e., if
		$$\operatorname{Vol}\left\{\theta \in U : \varphi^n(\theta) \in U,~\forall\,n \in \mathbb N \text{ and }\lim_{n \to \infty} \varphi^n(\theta) = \theta^*\right\}>0.$$
	\end{enumerate}
\end{definition}

\begin{proposition}\label{cekk_prop:LocGlobEqui}
	Let $\theta^*$ be a local minimum of $\mathcal L$ and suppose that the map $\varphi$ is non-singular. Then $\theta^*$ is locally Milnor stable if and only if it is Milnor stable.
\end{proposition}

\begin{proof}
	It follows directly from Definition \ref{cekk_def:Milnor} that every locally Milnor stable point is Milnor stable. It remains to show the converse. Suppose that $\varphi$ is non-singular and that $\theta^*$ is not locally Milnor stable and let $U\subset \R^D$ be a neighborhood of $\theta^*$ for which the set
	$$N\coloneqq \left\{\theta \in U : \varphi^n(\theta) \in U,~\forall\,n \in \mathbb N, \text{ and }\lim_{n \to \infty} \varphi^n(\theta) = \theta^*\right\}$$
	is a Lebesgue null set. By repeated application of the non-singularity of $\varphi$ we get that
	$$\operatorname{Vol}\left(\varphi^{-m}(N)\right) = 0, ~\forall\, m \in \mathbb N.$$
	By the definition of convergence, we have
	$$\left\{\theta \in \R^D : \lim_{n \to \infty} \varphi^n(\theta) = \theta^*\right\} = \bigcup_{m = 1}^\infty \varphi^{-m}(N)$$
	and thus 
	$$\operatorname{Vol}\left\{\theta \in \R^D : \lim_{n \to \infty} \varphi^n(\theta) = \theta^*\right\}=0$$
	implying that $\theta^*$ is not Milnor stable.
\end{proof}

\end{document}
